\chapter{Introduction}

\section{Motivation: From Theory to Practice}

This project investigates the security boundaries of \textbf{eBPF} (extended Berkeley Packet Filter) by moving from \textbf{theoretical vulnerability identification} to \textbf{practical exploitation}. The fundamental question is: \textit{Do logic bugs in the eBPF verifier actually allow memory corruption and privilege escalation in real kernel execution?}

A previous research team conducted a comprehensive, rule-based assessment of the XDP SynProxy program using ISO-IEC TS 17961-2013 as the vulnerability catalog. They identified approximately 60 potential violations and classified roughly 9 as theoretically exploitable by analyzing the verifier's behavior under code patches. Their work established the foundation for this study but left a critical gap: theoretical analysis alone is insufficient to understand real-world impact. This work bridges that gap by implementing functional proof-of-concept exploits, collecting multi-stage evidence, and determining which verifier bypasses represent genuine memory corruption primitives versus harmless runtime behaviors.

Building on the previous team's infrastructure and vulnerability catalog, we extended the testing environment with external-perspective exploitation capability, allowing attacks to be launched from the host machine against the target VM's network interface. This enhancement transforms the assessment from local-only testing into a realistic network-adjacent attack scenario. Additionally, we modified the VM provisioning scripts and network configuration to support shared file system access (via host-to-VM mount) and bridged networking, streamlining the exploit deployment and testing workflow.

The transition from theoretical to practical work requires:
\begin{enumerate}[label=\arabic*.]
  \item Engineering functional \textbf{Proofs of Concept (PoCs)} that trigger vulnerabilities through network packets
  \item Collecting \textbf{multi-stage evidence} (bytecode, verifier logs, translated code, runtime traces) to document the attack chain
  \item Distinguishing between \textit{runtime triggers} (which may be observable but harmless) and true \textbf{memory corruption primitives}  
  \item Documenting a repeatable, evidence-based methodology applicable to future eBPF verifier assessment
  \item Implementing an isolated, reproducible test environment with external-perspective exploitation capability
\end{enumerate}

\section{Scope and Objectives}

This work focuses on five selected vulnerabilities, chosen because they represent critical classes of verifier failures:

\begin{table}[h]
\centering
\begin{tabular}{|c|l|l|}
\hline
\textbf{ID} & \textbf{Vulnerability} & \textbf{Type} \\
\hline
5.06a & Function Pointer Type Mismatch & ABI/Argument confusion \\
5.06b & Wrong Number of Arguments & ABI/Argument confusion \\
5.06d & Wrong Argument Types & ABI/Argument confusion \\
5.39 & Unprototyped Function Pointer & Stack OOB write \\
5.46b & Tainted Length (Packet Buffer) & Heap/Stack OOB write \\
\hline
\end{tabular}
\caption{Selected vulnerabilities for practical exploitation study}
\end{table}

The objectives are to:
\begin{itemize}
  \item Demonstrate which verifier bypasses lead to observable runtime effects
  \item Quantify the severity of each vulnerability (ABI violation vs. memory corruption)
  \item Document evidence chains showing preservation through compilation, verification, and JIT
  \item Provide reproducible PoCs and analysis methodology
\end{itemize}

\section{Report Structure}

\begin{description}
  \item[Chapters 2--5] Establish eBPF background, threat model, methodology, and test environment
  \item[Chapter 6] Provide an overview of all five vulnerabilities and their classification
  \item[Chapters 7--11] Deep-dive analysis of each vulnerability (source, bytecode, verifier, runtime)
  \item[Chapters 12--15] Synthesize results, discuss limitations, propose future directions, and conclude
\end{description}
